\section*{Data and code availability}
The accompanying source package contains the English master, its Chinese translation, both supplements, the parameter record and the scripts used for the revised finite-element, matrix-scheme and controlled-discretization calculations. Running \texttt{python3 reproduction/run\_all.py} from the package root regenerates the mechanical baseline before the residual-control tests, related tracking extension, damping scan and framework-verification suite. Python~3, NumPy and SciPy are required for those calculations; Matplotlib is used to regenerate manuscript figures. The build instructions specify the separate PDF compilation steps.

The source manuscript is the \href{https://github.com/mop74720-commits/IJRR_CHNN_KAM_v10_20261007}{public v10 repository}, at commit \href{https://github.com/mop74720-commits/IJRR_CHNN_KAM_v10_20261007/tree/82519542dd057e8fbd6d6879ef33217761e8cbbb}{\texttt{82519542dd05}} (full identifier in the accompanying provenance record). Its S1/S2 archives remain the computational record of the original parameter set. They are not included as newly generated evidence, and the present revision does not claim to reproduce every historical result. The supplied Direction I research framework determines the three-subproblem organization of this revision. A source-to-section map in the package records where its scheme, controller, convergence arguments and numerical protocol enter the six-section article. The earlier migration table separately identifies which v10 results are retained as general theory, recomputed, restated conditionally or left with the historical benchmark.

\section*{Declarations}
This work reports mathematical analysis and computer-generated mechanical data. No human, animal, hardware or in-orbit experiment was performed for the revised benchmark. Material references and manufacturer specifications are cited as such; device calibration is not implied. This document is a research revision in the original article format and carries no publication or acceptance claim.
